\uuid{hXfB}
\exo7id{3409}
\titre{exo7 3409}
\auteur{quercia}
\organisation{exo7}
\datecreate{2010-03-10}
\isIndication{false}
\isCorrection{true}
\chapitre{Déterminant, système linéaire}
\sousChapitre{Système linéaire, rang}
\cpge{mpsi-mp2i.calcul-matriciel-systemes.c, pcsi.calcul-matriciel-systemes-lineaires.c, ptsi.calcul-matriciel-systemes-lineaires.c, mpsi-mp2i.determinants.d, pcsi.calcul-matriciel-systemes-lineaires.d, ptsi.calcul-matriciel-systemes-lineaires.d, mpsi-mp2i.matrices-applications-lineaires.c, pcsi.matrices-applications-lineaires.d, ptsi.matrices-applications-lineaires.d, bcpst1.systemes-lineaires.a, bcpst1.systemes-lineaires.c}
\module{Algèbre}
\niveau{L2}
\difficulte{}

\contenu{
\question{
\'Etudier l'existence de solutions du système :

 $$\left\{\begin{array}{ccccccc} x &-& my   &+& m^2z &=& m   \cr
           mx &-& m^2y &+& mz   &=& 1   \cr
           mx &+& y    &-& m^3z &=& -1. \cr\end{array}\right.$$
}
\reponse{
Système de Cramer ssi $m\ne 0,\pm1,\pm i$; compatible ssi $m\ne 0,\pm i$.
}
}
